\documentclass[10pt,a4paper]{amsart}
\usepackage[margin=19mm]{geometry}
\usepackage{amsmath,amssymb,mathtools}
\usepackage[scr=rsfs,cal=euler]{mathalfa}
\usepackage{xfrac}
\usepackage[unicode,hidelinks,bookmarks=false]{hyperref}
\newcommand{\ie}{i.e.\ }
\newcommand{\cf}{cf.\ }
\newcommand{\resp}{respectively\ }
\newcommand{\Subsection}{\subsection}
\providecommand{\nobreakdash}{\nobreak\hbox{-}}
\setlength{\emergencystretch}{2em}
\multlinegap0pt
\usepackage{bm}
\usepackage{bbm}
\usepackage{dsfont}
\usepackage{xspace}
\usepackage{cancel}
\usepackage{mathtools}
\usepackage{amsthm}
\makeatletter
\@ifundefined{theorem}{\newtheorem{theorem}{Theorem}}{}
\makeatother
\title[Fixed point results for contractions of polynomial type]{Fixed point results for contractions of polynomial type}
\author{Mohamed Jleli, Cristina Maria Pacurar, Bessem Samet}
\date{Source: arXiv:2406.03446v1 (CC BY 4.0)}
\begin{document}
\maketitle

\noindent\emph{Source and licence.} Fixed point results for contractions of polynomial type. Version arXiv:2406.03446v1 (CC BY 4.0), distributed under the Creative Commons Attribution 4.0 International licence, as recorded in the arXiv licence metadata for this exact version and shown on the abstract page https://arxiv.org/abs/2406.03446v1. Full rights notice: licenses/arxiv-2406.03446-CC-BY-4.0.txt.\par\smallskip

\noindent\textit{Editorial scope.} Counted unit: the Theorem whose source label is \texttt{T2.2} in the article (source0.tex lines 142--152, source label \texttt{T2.2}), delivered together with the article's own complete original proof of it (source0.tex lines 154--242, one \texttt{proof} environment of 89 lines). No mathematics is written, repaired or re-derived: statement and proof are the article's own text line for line. Required context retained uncounted: SPC (lines 132--134). Every definition, display and label of those lines is retained unchanged. Omitted from this package and not redistributed: 1--131, 135--141, 153--153, 243--913. Nothing inside a retained excerpt is omitted or altered: every retained body is delivered exactly as the article's lines are written, so no omission receipt of any kind accompanies this package. Citations: the retained excerpts carry no citation command at all, so no bibliography entry of the article is transcribed and no reference is redistributed. Reference adaptation: none was needed and none was made; every reference inside the delivered unit resolves inside it, so no unresolved reference or citation is produced. Printed slips kept literal: no printed word, symbol, hypothesis or number of the counted statement or of its proof was altered. Layout and class: the article's own class is \texttt{article} with packages including inputenc, xcolor, amssymb, amsmath, todonotes, amsthm, geometry, fancyhdr, upgreek, cite, enumitem, hyperref; the wrapper is the lane's pinned \texttt{amsart} editorial wrapper at 10pt on A4, which restates the theorem-like environments the retained text needs and adds the article's own macro definitions that the retained text needs (none), with extra wrapper packages (bm, bbm, dsfont, xspace, cancel, mathtools). The delivered numbering is the unit's own. Assets: no retained span contains a figure environment, a graphics-inclusion command, a PDF-page inclusion, a table or a TikZ picture; no image file is copied into this package, the package declares no asset dependency and no image file is redistributed with it. Licence: version arXiv:2406.03446v1 of this article is distributed under the Creative Commons Attribution 4.0 International licence, as recorded in the arXiv licence metadata for this exact version and shown on the abstract page https://arxiv.org/abs/2406.03446v1. A full rights notice is delivered at licenses/arxiv-2406.03446-CC-BY-4.0.txt. Characters: every retained excerpt is pure ASCII, so the delivered engine, XeLaTeX with its default Latin Modern text font, reports no missing character.
\begin{equation}\label{SPC}
\sum_{i=0}^ka_i(Tx,Ty)d^i(Tx,Ty)\leq \lambda \sum_{i=0}^k a_i(x,y)d^i(x,y)	
\end{equation}

\begin{theorem}\label{T2.2}
Let $(X,d)$ be a complete metric space and $T: X\to X$ be a polynomial contraction. Assume that the following conditions hold:
\begin{itemize}
\item[{\rm{(i)}}]$T$ is continuous;
\item[{\rm{(ii)}}] There exist $j\in\{1,\cdots,k\}$ and 	$A_j>0$ such that 
$$
a_j(x,y)\geq A_j,\quad x,y\in X.
$$
\end{itemize}
Then $T$ admits a unique fixed point $z^*\in X$. Moreover, for every $z_0\in X$, the Picard sequence  $\{z_n\}\subset X$ defined by $z_{n+1}=Tz_n$ for all $n\geq 0$, converges to $z^*$. 
\end{theorem}

\begin{proof}
We first prove that the set of fixed points of $T$ is nonempty.  Let $z_0\in X$ be fixed and  $\{z_n\}\subset X$ be the Picard sequence defined by  
$$
z_{n+1}=Tz_n,\quad  n\geq 0.
$$
Making use of \eqref{SPC} with $(x,y)=(z_0,z_1)$, we obtain
$$
\sum_{i=0}^ka_i(Tz_0,Tz_1)d^i(Tz_0,Tz_1)\leq \lambda \sum_{i=0}^k a_i(z_0,z_1)d^i(z_0,z_1),
$$
that is,
\begin{equation}\label{S1}
\sum_{i=0}^ka_i(z_1,z_2)d^i(z_1,z_2)\leq \lambda \sum_{i=0}^k a_i(z_0,z_1)d^i(z_0,z_1).
\end{equation}
Using again \eqref{SPC} with $(x,y)=(z_1,z_2)$, we obtain
$$
\sum_{i=0}^ka_i(Tz_1,Tz_2)d^i(Tz_1,Tz_2)\leq \lambda \sum_{i=0}^k a_i(z_1,z_2)d^i(z_1,z_2),
$$
that is,
$$
\sum_{i=0}^ka_i(z_2,z_3)d^i(z_2,z_3)\leq \lambda \sum_{i=0}^k a_i(z_1,z_2)d^i(z_1,z_2),
$$
which implies by \eqref{S1} that 
$$
\sum_{i=0}^ka_i(z_2,z_3)d^i(z_2,z_3)\leq \lambda^2 \sum_{i=0}^k a_i(z_0,z_1)d^i(z_0,z_1). 
$$
Continuing in the same way, we obtain by induction that   
\begin{equation}\label{S2}
\sum_{i=0}^ka_i(z_n,z_{n+1})d^i(z_n,z_{n+1})\leq \lambda^n \sum_{i=0}^k a_i(z_0,z_1)d^i(z_0,z_1),\quad n\geq 0.  	
\end{equation}
Since 
$$
a_j(z_n,z_{n+1})d^j(z_n,z_{n+1})\leq \sum_{i=0}^ka_i(z_n,z_{n+1})d^i(z_n,z_{n+1}),
$$
we obtain by (ii) that 
$$
A_j d^j(z_n,z_{n+1})\leq \sum_{i=0}^ka_i(z_n,z_{n+1})d^i(z_n,z_{n+1}),
$$
which implies by \eqref{S2} that 
\begin{equation}\label{S3}
d^j(z_n,z_{n+1})\leq  \lambda^n \sigma_{j,0},\quad n\geq 0,	
\end{equation}
where 
\begin{equation}\label{sigma0}
\sigma_{j,0}=A_j^{-1}\sum_{i=0}^k a_i(z_0,z_1)d^i(z_0,z_1).
\end{equation}
Then, making use of \eqref{S3} and the triangle inequality, we obtain that for all $n\geq 0$ and $m\geq 1$,
$$
\begin{aligned}
d^j(z_n,z_{n+m})&\leq d^j(z_n,z_{n+1})+d^j(z_{n+1},z_{n+2})+\cdots+d^j(z_{n+m-1},z_{n+m})\\
&\leq \sigma_{j,0}\left(	\lambda^n+\lambda^{n+1}+\cdots+\lambda^{n+m-1}\right)\\
&=\sigma_{j,0} \lambda^n \frac{1-\lambda^{m}}{1-\lambda}\\
&\leq \sigma_{j,0}  \frac{\lambda^n}{1-\lambda},
\end{aligned}
$$
which yields
$$
d(z_n,z_{n+m})\leq \left(\frac{\sigma_{j,0} }{1-\lambda}\right)^{\frac{1}{j}} (\lambda^{\frac{1}{j}})^n\to 0\mbox{ as }n,m\to \infty.
$$
This shows that $\{z_n\}$ is a Cauchy sequence. Since $(X,d)$ is complete, there exists $z^*\in X$ such that 
$$
\lim_{n\to \infty}d(z_n,z^*)=0,
$$
which implies by the continuity of $T$ that 
$$
\lim_{n\to \infty}d(z_{n+1},Tz^*)=\lim_{n\to \infty}d(Tz_n,Tz^*)=0.
$$
Then, from the uniqueness of the limit, we deduce that $Tz^*=z^*$, that is, $z^*$ is a fixed point of $T$. 

We now show that $z^*$ is the unique fixed point of $T$. Indeed, if $z^{**}\in X$ is another fixed point of $T$, i.e., $Tz^{**}=z^{**}$ and $d(z^*,z^{**})>0$, then making use of \eqref{SPC} with $(x,y)=(z^*,z^{**})$, we get 
$$
\sum_{i=0}^ka_i(Tz^*,Tz^{**})d^i(Tz^*,Tz^{**})\leq \lambda \sum_{i=0}^k a_i(z^*,z^{**})d^i(z^*,z^{**}),	
$$ 
that is,
\begin{equation}\label{S4}
\sum_{i=0}^k a_i(z^*,z^{**})d^i(z^*,z^{**})\leq \lambda \sum_{i=0}^k a_i(z^*,z^{**})d^i(z^*,z^{**}).
\end{equation}
On the other hand, from	(ii), we have 
$$
\begin{aligned}
\sum_{i=0}^k a_i(z^*,z^{**})d^i(z^*,z^{**}) &\geq  a_j(z^*,z^{**})d^j(z^*,z^{**})\\
&\geq A_j d^j(z^*,z^{**}).
\end{aligned}
$$
Since $A_j>0$ and $d(z^*,z^{**})>0$, we deduce that 
$$
\sum_{i=0}^k a_i(z^*,z^{**})d^i(z^*,z^{**})>0. 
$$
Then, dividing  \eqref{S4} by $\sum_{i=0}^k a_i(z^*,z^{**})d^i(z^*,z^{**})$, we reach a contradiction with $\lambda\in (0,1)$. Consequently, $z^*$ is the unique fixed point of $T$. This completes the proof of Theorem \ref{T2.2}. 
\end{proof}

\end{document}
